\documentclass[11pt]{article}
\usepackage[margin=1in,top=0.8in]{geometry}
\usepackage[T1]{fontenc}
\usepackage{libertine}
\usepackage{microtype}
\usepackage{amsmath}
\usepackage{amssymb}
\usepackage{xcolor}
\usepackage{listings}
\usepackage{enumitem}
\usepackage{titlesec}
\usepackage{hyperref}

% Prevent widows and orphans
\widowpenalty=10000
\clubpenalty=10000
\raggedbottom

% Define colors for code
\definecolor{codegray}{rgb}{0.5,0.5,0.5}
\definecolor{codeblue}{rgb}{0.1,0.1,0.5}
\definecolor{backcolour}{rgb}{0.97,0.97,0.97}
% Code listing style
\lstdefinestyle{mystyle}{
	backgroundcolor=\color{backcolour},   
	basicstyle=\footnotesize\ttfamily,
	breakatwhitespace=false,         
	breaklines=true,                 
	captionpos=b,                    
	keepspaces=true,                 
	showspaces=false,                
	showstringspaces=false,
	showtabs=false,                  
	tabsize=2,
	frame=single,
	framerule=0.5pt,
	rulecolor=\color{gray!30},
	xleftmargin=0.5em,
	xrightmargin=0.5em,
	framexleftmargin=0.5em,
	framexrightmargin=0.5em,
	aboveskip=0.5em,
	belowskip=0.5em
}
\lstset{style=mystyle}
% Title formatting
\title{\Large\textbf{What If We Trained LLMs on Thoughts About Text?}}
\author{Henrik Westerberg\\
	\small{henrik.westerberg@emergentwisdom.org}}
\date{\today}

\begin{document}
	
	\maketitle
	
	\begin{abstract}
		\noindent
		Current large language models excel at generating human-like text but evaluate ideas only implicitly through their training data. We propose \textit{commentary-augmented training}: teaching LLMs to predict text and parallel commentary simultaneously, explaining why that text is valuable, flawed, or profound. Unlike methods that add evaluation post-hoc, our approach creates computational entanglement—making the LLM's ``inner monologue'' visible during generation. While the scale required necessitates AI-generated commentary, we hypothesize that this joint prediction fundamentally changes how models learn. This could enable: (1) real-time evaluation during generation; (2) transparent reasoning in outputs; (3) integrated ethical alignment; and (4) self-improvement through iterative refinement. Teaching machines to predict not just what we say but what we think about what we say represents a promising direction for AI systems that can explicitly reason about the quality and implications of their outputs.
	\end{abstract}
	
	\vspace{0.5em}
	
	\section{The Idea}
	
	Large language models have demonstrated remarkable capabilities in generating human-like text, yet they learn to evaluate ideas only implicitly through the patterns in their training data. While these models can explain why E=mc² is considered profound, they do so by pattern-matching human explanations rather than through explicit evaluative reasoning. They excel at predicting what comes next in a sequence, but lack a dedicated mechanism for assessing what truly matters \cite{brown2020language}.
	
	The problem runs deeper: the structure of human knowledge---papers citing papers, links between ideas---imperfectly reflects epistemic value. Well-cited ideas aren't always true; profound insights sometimes go unrecognized for decades. We often evaluate based on frequency, authority, and social consensus rather than systematic assessment of intrinsic value. Current LLMs absorb these biases, learning to reproduce human judgment patterns without questioning their validity or developing alternative evaluation criteria.
	
	We propose rethinking LLM training: teach models not merely on text, but on parallel commentary that explicates why that text is valuable, flawed, or profound. Rather than learning from raw text alone, models would learn from text paired with sophisticated analysis of that text. It's the difference between training a student by having them read textbooks versus training them with textbooks annotated by experts who explain why each concept matters, what mistakes to avoid, and how ideas connect.
	
	This commentary captures what the web of references leaves implicit---the actual reasoning about epistemic value. By learning to predict both the primary content and its evaluative commentary simultaneously, LLMs can develop explicit meta-cognitive abilities. In essence, we make the LLM's ``inner monologue'' visible---revealing not just what the model generates but its real-time assessment of that content. We hypothesize that training on parallel commentary will lead to parameters that simultaneously encode both what to say and how to evaluate what is said---creating representations where generation and evaluation are intrinsically linked rather than evaluation emerging only as an implicit capability.
	
	This linkage occurs through computational entanglement. When a model learns to predict both text and commentary as parallel tracks, its hidden representations must encode both capabilities simultaneously. The model cannot access its knowledge about `Einstein' without also accessing evaluative concepts like `revolutionary', `testable', and `unified theory'---these become inseparable in the learned representations. Unlike post-hoc evaluation that requires a separate forward pass after generation, this entanglement means evaluation is woven into the representational fabric that produces each token. The model doesn't first generate then evaluate; it generates through evaluation.
	
	The transformative potential lies in how this approach compounds: by training on billions of evaluations, the model learns not just individual judgments but the underlying patterns of evaluation itself. An LLM that internalizes how humans evaluate---not just what they conclude---becomes capable of applying sophisticated judgment to entirely novel domains and ideas.
	
	\section{The Method}
	
	Traditional LLM training follows a straightforward next-token prediction paradigm:
	\begin{lstlisting}
		
		"The equations of general relativity show that..." -> [predict next token]
	\end{lstlisting}
	Our proposed commentary-augmented training introduces a parallel evaluative track:
	\begin{lstlisting}
		
		[TEXT]: "The equations of general relativity show that spacetime curves in the presence of mass-energy"
		[COMMENTARY]: "Accurate summary of Einstein's 1915 breakthrough. Core insight: geometry equals gravity. Confirmed via light bending. Unifies separate concepts."
		
		[TEXT]: "Consciousness creates gravity through quantum vibrations"
		[COMMENTARY]: "Novel connection attempt. Lacks mechanism. No mathematical framework. Must address quantum decoherence. Needs testable predictions."
	\end{lstlisting}
	In this approach, the LLM would learn to predict both tracks simultaneously, internalizing not merely patterns of language but patterns of evaluation. When the model sees thousands of examples where 'Confirmed via light bending' appears for testable predictions and 'Lacks mechanism' appears for vague claims, it would learn to recognize these deeper patterns. The commentary makes explicit what traditional training leaves implicit: why Einstein's geometric insight warranted belief: 'empirical predictions', 'mathematical framework' while consciousness-gravity claims do not: 'no mechanism', 'scale mismatch'. Through this systematic exposure, the model would develop criteria for distinguishing profound insights from superficial connections.
	
	These commentaries would need to capture how evaluation varies by discipline and framework. Consider how even within physics, researchers reach opposite conclusions about string theory---mathematical physicists might praise its ``profound unification of quantum mechanics and general relativity'' while experimentalists criticize it as ``mathematically impressive but empirically sterile.'' By training on these conflicting expert assessments, the model would learn that evaluation depends on which criteria matter most: empirical testability, mathematical coherence, explanatory power, or practical utility. Through exposure to these legitimate disagreements, the model learns that rigorous evaluation can lead to different conclusions depending on one's epistemic values.
	
	Crucially, commentary should depend on context---claims appearing without proper justification receive more scrutiny than those following systematic argumentation. A statement like ``this treatment reduces symptoms by 40\%'' might receive commentary noting 'lacks control group, sample size, and statistical significance' when appearing alone, but 'well-supported conclusion given the preceding double-blind study' when following proper methodology. The model would learn that strong evaluation considers not just individual statements but their epistemic support: whether evidence precedes assertions, whether logical connections are explicit, and whether appropriate uncertainty is acknowledged. This teaches intellectual rigor at the structural level, not just the statement level.
	
	Commentary must operate at multiple scales. Word-level evaluation---'precise' vs 'vague'---teaches linguistic discrimination. Sentence-level commentary identifies logical flaws and unsupported claims. Paragraph-level assessment captures argumentative coherence. Document-level evaluation judges whether conclusions follow from premises. Training on mixed granularities would teach models to evaluate at the appropriate level---catching both local errors and global inconsistencies.
	
	\section{The Data Generation Challenge}
	
	The most immediate obstacle facing commentary-augmented training is the sheer scale of annotation required. Training modern LLMs requires trillions of tokens; creating human-written commentary for even a fraction of this data would be prohibitively expensive and logistically impossible. The only practical path forward is to leverage existing LLMs to generate the bulk of evaluative commentary.
	
	This apparent circularity---using LLMs to create training data for better LLMs---might seem to undermine our approach. If current models lack explicit evaluative reasoning, how can they generate the commentary needed to teach it? However, this critique misunderstands what we're proposing. When prompted appropriately, current LLMs can already produce sophisticated evaluative language by recombining patterns from their training. The key insight is that making this evaluative reasoning explicit in the training data---even if initially generated through prompted pattern-matching---fundamentally changes how future models learn to process information.
	
	Consider an analogy: a student who can solve physics problems intuitively might struggle to explain their reasoning. But when forced to write out their thought process step-by-step, they often discover gaps in understanding and develop more rigorous thinking. The same principle applies when generating training data---the dialogue between our evaluative prompts and the model's responses creates something new: structured commentary that emerges from their interaction. This prompt-response activity doesn't just extract hidden knowledge but generates newly organized information that future models can internalize as explicit evaluative capabilities.
	
	\textbf{The Bootstrapping Advantage}: Rather than a weakness, AI-generated commentary offers several advantages: (1) systematic coverage---unlike human annotators who might focus on interesting cases, AI can methodically generate commentary for every example, ensuring no capability appears without evaluation; (2) controlled variation---by systematically varying prompts (e.g., "evaluate from a safety perspective," "consider long-term consequences," "assess epistemic rigor"), we can generate genuinely different evaluative lenses for the same content; (3) iterative refinement---as each generation develops stronger evaluation capabilities through training, it can generate more nuanced commentary for the next generation's training data; and (4) cost-effectiveness---the only approach economically viable at the required scale.
	
	The quality bar isn't perfection---it's whether the commentary makes evaluative reasoning more explicit than current training data. Even commentary that ``merely'' reorganizes implicit knowledge into explicit evaluative statements provides value. When a model trains on text paired with commentary explaining 'this lacks empirical support' or 'this reasoning contains a hidden assumption', it learns to recognize these patterns more directly than inferring them from raw text alone.
	
	The ultimate test remains empirical: do models trained on AI-generated commentary demonstrate improved evaluative capabilities compared to standard training? We hypothesize yes, because explicit evaluative language---regardless of its origin---provides clearer learning signals than leaving evaluation entirely implicit.
	
	\section{The Implementation}
	
	With these challenges in mind, we now turn to the practical implementation of commentary-augmented training. Commentary-augmented training uses standard transformer architecture with modified training data. Several architectural approaches merit consideration: separate output heads (clean task separation but less integration), dual networks (specialized optimization), or fully integrated shared parameters (maximum cross-pollination but mode-switching challenges). We suggest starting with fully integrated parameters to maximize interaction between generation and evaluation capabilities.
	
	The training objective combines both tasks: $\mathcal{L} = \mathcal{L}_{\text{text}} + \alpha \mathcal{L}_{\text{commentary}}$, where $\alpha$ weights the importance of commentary prediction. While optimal tuning remains empirical, starting with equal weighting ($\alpha = 1$) signals that evaluation matters as much as generation. We hypothesize that data quality matters more than architectural complexity---that a standard transformer trained on rich text-commentary pairs would outperform sophisticated architectures trained on text alone.
	
	Testing focuses on two hypotheses. First, that commentary-trained models show improved evaluation capabilities on tasks such as argument flaw detection, claim verification, and generating critiques of held-out texts. These models should express appropriate uncertainty at knowledge boundaries and provide coherent self-evaluations when prompted. Second, that this training enhances standard generation through increased claim support rates, logical consistency, and error self-correction. The first hypothesis has strong face validity; the second remains genuinely uncertain. These capabilities might develop independently---improved evaluation without improved generation would still validate the approach while suggesting architectural refinements.

	\section{Implications}
	
	\textbf{Enhanced Evaluation}: An LLM trained through this methodology would develop evaluative capabilities that operate during generation, not merely after. The model could assess any content---user proposals, external arguments, or its own reasoning---in real-time. When asked ``What are the flaws in this argument?'' or ``Is this approach sound?'', it would provide sophisticated analysis grounded in learned evaluative patterns. This reflexive capability would serve dual purposes: catching logical errors and recognizing exceptional insights. By learning from commentary that explicitly marks strong reasoning, identifies logical flaws, and delineates knowledge boundaries, the model could develop nuanced judgment about when ideas are well-supported, when connections are meaningful, and when claims exceed available evidence.
	
	\textbf{Transparent Reasoning}: Unlike systems that generate outputs as black boxes, these models are designed to surface evaluative thinking as a natural part of their responses. Users would gain insight into what factors the model considers, what uncertainties it recognizes, and what trade-offs it weighs---not through special prompting but because the model learns through training that good responses include metacognitive transparency. While this wouldn't solve mechanistic interpretability, it would provide something immediately useful: models that acknowledge their own limitations, express appropriate uncertainty, and explain their confidence levels. This evaluative transparency would help users calibrate when to trust outputs and understand the reasoning patterns shaping the model's responses.
	
	\textbf{Alignment}: Every capability the model acquires comes paired with its ethical dimensions in the training data. When learning about persuasion, it simultaneously encounters commentary about manipulation; when understanding chemistry, it sees dual-use concerns. Values aren't imposed through external constraints but integrated into how the model learns to evaluate everything it knows. Research on multi-task learning shows that when models optimize for correlated tasks using shared parameters, they develop unified representations that encode both aspects efficiently \cite{ruder2017overview}. Through this joint optimization, the system learns to treat ``how to do X'' and ``when X might be harmful'' as inseparable concepts---ethical reasoning becomes part of its core evaluative process through the commentary-augmented training.
	
	\textbf{Self-Improvement Loop}: This approach could enable bootstrapping if two conditions hold: (1) training on commentary improves evaluation ability, and (2) better evaluators generate better commentary. When both conditions are met, the model could enhance its own training data---each iteration potentially producing richer commentary than the last. The model would become a progressively better teacher of itself, creating new evaluative insights rather than just filtering existing outputs. Each cycle could compound: better evaluation $\rightarrow$ better commentary $\rightarrow$ better next-generation model. Of course, this assumes quality improvements rather than degradation---a critical empirical question we discuss in the following section.
	
	\section{Limitations}
	
	We present commentary-augmented training not as a complete solution to machine understanding, but as a promising direction that could enhance current capabilities. Whether this approach enables ``genuine'' understanding remains an open philosophical question---no more or less mysterious than the same question applied to standard LLMs. What we propose is empirically testable: can models trained with evaluative commentary demonstrate improved judgment on novel instances?
	
	A fundamental challenge lies in distinguishing genuine evaluative capability from sophisticated mimicry. Models might learn to generate plausible evaluative language without deeper understanding, creating an ``evaluative uncanny valley''---superficially sophisticated commentary that misses crucial nuances human experts would catch. Additionally, available commentary data represents polished post-hoc analysis rather than real-time cognitive processes---the gap between formal commentary and actual thinking may limit the approach's transformative potential while still enabling meaningful improvements.
	
	The approach also faces a deeper philosophical challenge: which evaluative frameworks should we include? As Hume observed, we cannot derive ``ought'' from ``is''---no amount of factual training data can definitively establish which values or evaluative criteria are correct. The selection of commentary inevitably embeds particular philosophical and cultural assumptions about what constitutes good reasoning, valid evidence, or ethical behavior.
	
	This limitation, however, highlights why commentary-augmented training matters for AI safety. Given our limited ability to understand or control how models process information internally, models will encounter texts like Dante's Inferno regardless. Without commentary, they learn descriptions of torture as neutral text patterns. With commentary—ranging from fundamental moral grounding: 'this depicts intense physical suffering, which is inherently harmful' to contextual analysis: 'medieval allegory reflecting period theology' to critical perspective: 'artistic achievement that modern readers approach with ethical reservations'—models learn to process violent content within evaluative frameworks rather than as value-neutral data.
	
	While we cannot guarantee models will be constrained by these moral judgments, the choice is clear. Authors throughout history have written for readers with bodies that feel pain, cultural contexts that shape meaning, and innate empathic responses. AI systems possess none of these. Commentary-augmented training attempts to bridge this gap, making explicit the ethical understanding that human readers bring implicitly to any text. We developed this approach with careful consideration of AI safety concerns, believing that transparent evaluative reasoning may be essential for maintaining human oversight as AI systems grow more capable.
	
	The self-improvement loop presents both opportunities and risks. Like all iterative training approaches, commentary enhancement faces potential degeneracy challenges---if initial commentary quality is poor or biased, self-improvement could amplify rather than correct these issues. A system recursively enhancing its evaluative capabilities could develop in unexpected ways: it may experience model collapse, where training on self-generated data causes ``irreversible defects'' and leads to outputs the original model would never produce \cite{shumailov2023curse}, or engage in reward hacking, exploiting flaws in its evaluative signals rather than learning intended behaviors \cite{gao2025reward}. These specific mechanisms exemplify broader concerns about recursive self-improvement in AI systems, where enhanced capabilities can lead to unforeseen and potentially uncontrollable development trajectories \cite{bostrom2014superintelligence}.
	
	However, unlike opaque self-improvement, commentary-augmented models would document their reasoning at each step, providing visibility into their evolving judgment. This transparency, combined with standard safeguards like diverse data sources and held-out validation, offers paths to managing these risks while potentially creating natural alignment through integrated ethical reasoning---though this requires careful empirical validation.
	
	\section{Related Work}
	
	Commentary-augmented training represents a fundamental departure from existing approaches to improving LLM capabilities. While current methods focus on shaping model behavior through external mechanisms---rewards, instructions, or constraints---our approach integrates evaluative reasoning directly into the learning process itself. Rather than teaching models what to output or how to behave, we teach them to predict the evaluative thinking that accompanies all meaningful text.
	
	\textbf{Reinforcement Learning from Human Feedback (RLHF)} \cite{christiano2017deep} optimizes outputs according to human preferences post-generation, shaping what models choose to output without altering how they represent quality during generation. While RLHF learns what humans prefer, our approach teaches models to predict and generate the reasoning behind those preferences. Crucially, where RLHF externalizes evaluation into a separate reward model, we propose making evaluation intrinsic to the generation process itself.
	
	\textbf{Instruction tuning} \cite{wei2021finetuned} can produce models that analyze text when asked, but treats evaluation as just another text generation task. Our approach aims to make evaluation a parallel cognitive process---the model learns to evaluate while generating, not after being asked. The distinction is between learning to perform evaluation as a task versus learning to think evaluatively as an inherent aspect of language processing.
	
	\textbf{Chain-of-thought prompting} \cite{wei2022chain} elicits reasoning from models trained on standard text. While effective, it reveals capabilities the model already has. Commentary-augmented training builds evaluation directly into the training process, rather than hoping it emerges implicitly. Moreover, while CoT focuses on procedural reasoning steps explaining \emph{how} to derive an answer, our commentary explicitly addresses \emph{why} text is valuable, flawed, or profound---a fundamentally evaluative rather than procedural focus.
	
	\textbf{Constitutional AI} \cite{bai2022constitutional} implements behavioral constraints through rules, focusing on adherence rather than developing evaluative capabilities. Where Constitutional AI says ``don't generate harmful content,'' our method teaches models to recognize and articulate why content might be harmful through explicit commentary during training.
	
	Recent work increasingly leverages AI-generated evaluative data for training. \textbf{Critique Fine-Tuning} \cite{wang2025critiquefinetuninglearningcritique} trains models to critique incorrect responses using AI-generated critique datasets, achieving 4-10\% performance improvements despite ~20\% error rates in the generated critiques. \textbf{CREST} \cite{lee2025selftrainingmeetsconsistencyimproving} uses self-generated rationales filtered by consistency checks for self-training in QA tasks. \textbf{CTRL} \cite{xie2025teachinglanguagemodelscritique} trains separate critic models using AI-synthesized feedback guided by execution results. These successes demonstrate that making implicit evaluative knowledge explicit through AI generation---even when imperfect---provides valuable training signals.
	
	However, these methods still treat evaluation as a specialized capability learned separately from generation. CFT teaches models to produce critiques of existing text (learning $P(\text{critique}|\text{text})$), while our approach trains models to generate text and commentary simultaneously (learning $P(\text{text}, \text{commentary}|\text{context})$). This architectural difference is fundamental: rather than adding evaluative capabilities through post-hoc training or separate models, we propose making evaluation and generation computationally entangled from the outset, forcing models to develop representations where generating text and evaluating it are inseparable aspects of the same process.
	
	\section{Conclusion}
	
	We asked: what if we trained LLMs on thoughts about text? This approach aims to provide them with something humans often take for granted: constant evaluative judgment. Every piece of human text carries an ``invisible shadow''---the evaluative thinking that shaped its creation. Writers constantly assess: `this needs evidence', `this connection matters', `this word is precise'. Readers bring their own judgments: `this lacks rigor', `this argument has gaps', `this insight is profound'. Current LLMs learn from our textual output but can only infer these hidden layers of metacognition indirectly, piecing together what makes text valuable from patterns rather than accessing the reasoning itself.
	
	The timing for exploring commentary-augmented training appears particularly relevant. As LLMs become more powerful and pervasive, their lack of explicit evaluative reasoning raises important concerns. Models can generate plausible-sounding misinformation, struggle to acknowledge uncertainty appropriately, and cannot explain their confidence levels. Meanwhile, recent work demonstrates that AI-generated evaluative data can improve model capabilities---CFT achieved 4-10\% improvements despite imperfect critiques, while CREST successfully leveraged self-generated rationales. These successes suggest the infrastructure and precedent exist for more ambitious integration of evaluation into generation.
	
	We have presented commentary-augmented training: teaching models to predict text and commentary as parallel, entangled tracks. This differs from multi-task learning or post-hoc evaluation by proposing that models learn P(text, commentary|context) rather than P(text|context) followed by P(evaluation|text). The hypothesis is that this creates computational entanglement where the model cannot access knowledge without simultaneously accessing evaluative concepts. Making the LLM's ``inner monologue'' visible could emerge not as an add-on feature but as an inherent aspect of how it processes language.
	
	The potential implications present both opportunities and risks. A model that generates through evaluation might better recognize uncertainties and provide transparent reasoning---valuable for AI safety. Yet this same capability could accelerate the development of more powerful systems. The bootstrapping possibility, where each generation creates better commentary for training the next, exemplifies this duality: it promises iterative improvements in metacognitive capabilities while raising concerns about recursive self-improvement that have long worried safety researchers. However, the explicit nature of commentary might offer a unique advantage---allowing us to monitor capability development through the very transparency it creates.
	
	Empirical validation remains essential. Initial experiments could implement commentary-augmented training on existing architectures and measure whether models demonstrate enhanced evaluation on held-out tasks. Success indicators might include spontaneous uncertainty expression, self-correction during generation, and coherent self-evaluation when prompted. The recent advances in AI-generated evaluative data across multiple paradigms suggest promising foundations for such experiments.
	
	So, what if we trained LLMs on thoughts about text? We might develop machines that don't merely generate human-like text but begin to engage in the evaluative thinking that makes such text worth generating. As AI systems increasingly shape human knowledge and decision-making, teaching them not just what we say but what we think about what we say represents an important research direction---one that could help bridge the gap between pattern matching and genuine understanding, even as we remain mindful of the capabilities such understanding might unlock.
	
	\clearpage
	\bibliographystyle{unsrt}
	\bibliography{references}
	
\end{document}